% Subsection 2.2.1: 全量微调的局限性

全量微调是模型适配下游任务最直接的方式，其通过更新预训练模型的全部参数，使模型整体向目标任务分布靠拢。但在算法代码生成这类场景下，该方式存在明显短板：
\begin{itemize}
    \item \textbf{资源开销高}：现代代码大模型参数规模普遍达到数十亿量级，全量微调对显存、算力与训练时间要求极高，普通研究环境难以支撑稳定收敛与多次对照实验；
    \item \textbf{易过拟合}：在小规模专用数据集上，全量微调更容易陷入过拟合，导致模型在标准测试集上泛化能力下降；
    \item \textbf{存储与切换成本大}：每次微调都会产生完整的权重副本，不利于多方案、多配置的快速对比。
\end{itemize}
因此，面向算法代码生成任务开展研究时，需要更轻量、更稳定的参数高效微调方案。